\chapter{Technology-Independent Optimization}
\label{ch:opt}

This chapter describes the AIG restructuring passes that shrink area (AND-node
count) and depth (levels) \emph{before} any technology mapping. Every pass shares
the same shape: a lightweight \emph{driver} walks the network node-by-node, an
\emph{engine} proposes a local replacement, and a \emph{gain rule} accepts the
move only when it pays off. The passes differ in how they enumerate the local
neighborhood (cuts, cones, supergates) and how they synthesize a candidate
(precomputed subgraphs, factoring, resubstitution, or SAT).

\glance{Balance retimes AND-trees for depth; \cmd{rewrite}/\cmd{drw} swap in
precomputed 4-input subgraphs; \cmd{refactor} re-factors a reconvergent cone;
\cmd{resub} reconnects a node to existing divisors; \cmd{mfs} uses SAT and
don't-cares. All are DAG-aware: a move is kept only if the MFFC freed exceeds the
new nodes added (hash-shared nodes counted as free).}

\section{Balance}
\label{sec:opt-balance}

Balancing does not change functionality or (much) area; it re-associates
multi-input AND/EXOR trees so that late-arriving inputs sit near the root,
minimizing depth. Two independent implementations exist.

\paragraph{Classic AIG (\cmd{balance}).}
The driver \id{Abc\_NtkBalance} (\file{src/base/abci/abcBalance.c}) rebuilds the
network from the COs down. For each node, \id{Abc\_NodeBalanceCone} collects the
maximal ``implication supergate'' --- the transitive AND-cone that stops at a
complemented edge, a PI, or (when neither \id{fDuplicate} nor \id{fSelective} is
set) a node with more than one fanout; the recursion also hard-caps the supergate
at $10000$ literals. \id{Abc\_NodeBalance\_rec} then sorts the leaves by
decreasing level and repeatedly pops the two lowest-level operands and ANDs them,
producing a near-balanced tree. Two heuristics inject a \emph{sharing bias}:

\begin{itemize}
\item \id{Abc\_NodeBalanceFindLeft} finds the left-most operand of the same level
as the last one, defining the legal range of partners (only when
\id{fUpdateLevel} is on; otherwise the bound is $0$).
\item \id{Abc\_NodeBalancePermute} scans that range and, via
\id{Abc\_AigAndLookup}, prefers a partner whose AND already exists in the
structural hash, reusing a node instead of creating one.
\end{itemize}

\paragraph{GIA (\cmd{\&b}).}
\id{Gia\_ManBalance} (\file{src/aig/gia/giaBalAig.c}) is the scalable twin used in
\cmd{\&} recipes. \id{Gia\_ManSuperCollect} gathers AND and XOR supergates,
capping each at $50$ collected literals (\id{Vec\_IntSize(p->vSuper) > 50}), then
\id{Gia\_ManSimplifyAnd}/\id{Gia\_ManSimplifyXor} fold constants and duplicates.
\id{Gia\_ManBalanceGate} sorts operands by level and combines the two lowest.
To bias toward sharing without quadratic cost, \id{Gia\_ManPrepareLastTwo} only
inspects a \emph{lookback window of the last 8 operands}
(\id{Stop = Abc\_MaxInt(Stop, nSize - 9)}) when hunting for a hash-shareable pair.

\paragraph{Area balancing (\cmd{\&b -a}).}
\id{Gia\_ManAreaBalance} $\rightarrow$ \id{Dam\_ManAreaBalanceInt} runs a global
divisor-extraction loop (\id{Dam\_Man\_t}) that repeatedly pulls the most frequent
two-input pair out of all supergates. A priority queue ranks divisors by
occurrence count plus a small tie-breaker $0.005 \times$ slack, where the slack
computed by \id{Dam\_ManDivSlack} is clamped to $100$. The loop stops after
\id{nNewNodesMax} extractions or once the top priority drops below $2$.

\section{Rewrite (\texorpdfstring{\cmd{rewrite}}{rewrite})}
\label{sec:opt-rewrite}

DAG-aware rewriting replaces 4-input cones with equivalent, smaller precomputed
subgraphs. Driver \id{Abc\_NtkRewrite} (\file{src/base/abci/abcRewrite.c}) sets up
a cut manager and visits each node once; engine \id{Rwr\_NodeRewrite}
(\file{src/opt/rwr/rwrEva.c}) does the work.

\paragraph{Cuts.}
\id{Abc\_NtkStartCutManForRewrite} fixes the cut parameters
(\file{src/base/abci/abcRewrite.c}):

\begin{center}
\begin{tabular}{@{}ll@{}}
\toprule
\textbf{Parameter} & \textbf{Value} \\
\midrule
\id{nVarsMax} (cut size $k$)          & 4 \\
\id{nKeepMax} (max cuts kept / node)  & 250 \\
\id{fTruth} (compute truth tables)    & 1 \\
\id{fFilter} (filter dominated cuts)  & 1 \\
\bottomrule
\end{tabular}
\end{center}

\id{Rwr\_NodeRewrite} considers \emph{only} the 4-input cuts
(\id{pCut->nLeaves < 4} are skipped). The 16-bit truth table of each cut is
NPN-canonized through precomputed permutation/phase tables (\id{p->pPerms},
\id{p->pPhases}) into one of the \textbf{222} NPN classes of 4-variable functions
(\id{p->vClasses->nSize == 222}; the function table \id{pTable} spans
\id{nFuncs = 1<<16 = 65536} entries, set in \file{src/opt/rwr/rwrMan.c}). Each
class points to a small forest of precomputed replacement subgraphs
(\id{Rwr\_Node\_t}, decomposition graphs).

\paragraph{Gain rule.}
For each candidate the MFFC of the root is labeled
(\id{Abc\_NodeMffcLabelAig}, giving \id{nNodesSaved}) and \id{Rwr\_CutEvaluate}
computes, per subgraph, how many \emph{new} nodes would be added via
\id{Dec\_GraphToNetworkCount} --- crucially, nodes that already exist in the
structural hash are reused and counted as free. The gain is

\[
\text{gain} = \text{nNodesSaved (MFFC freed)} - \text{nNodesAdded (new nodes)} .
\]

The move is accepted when \id{nGain > 0}, or when \id{nGain == 0} and zero-cost
replacements are enabled (\cmd{rewrite -z}, the \id{fUseZeros} flag). Nodes that
are persistent or have more than $1000$ fanouts are skipped; a cut is also
dropped if more than two of its leaves are singly-fanout nodes.

\section{DAR rewriting (\texorpdfstring{\cmd{drw} / \cmd{dc2}}{drw / dc2})}
\label{sec:opt-dar}

\id{Dar\_ManRewrite} (\file{src/opt/dar/darCore.c}) is the AIG-package twin of
\cmd{rewrite} used inside \cmd{dc2} and exposed as \cmd{drw}. It is faster because
it enumerates cuts on the fly and stores its own subgraph library
(\id{Dar\_LibPrepare}). Defaults come from \id{Dar\_ManDefaultRwrParams}:

\begin{center}
\begin{tabular}{@{}lll@{}}
\toprule
\textbf{Field} & \textbf{Default} & \textbf{Meaning} \\
\midrule
\id{nCutsMax}  & 8 & max cuts computed per node \\
\id{nSubgMax}  & 5 & max subgraphs tried (a ``magic number'') \\
\id{nMinSaved} & 1 & min nodes saved to accept \\
\id{fUpdateLevel} & 0 & keep depth bounded by required time \\
\id{fUseZeros} & 0 & allow zero-gain moves (sets \id{nMinSaved=0}) \\
\bottomrule
\end{tabular}
\end{center}

Three-leaf cuts are padded to four (\id{pCut->pLeaves[nLeaves++] = 0}) before
\id{Dar\_LibEval}. A replacement is accepted only when
\id{p->GainBest >= nMinSaved}; \id{Dar\_LibBuildBest} then constructs it and
\id{Aig\_ObjReplace} splices it in.

\section{Refactor (\texorpdfstring{\cmd{refactor}}{refactor})}
\label{sec:opt-refactor}

Refactoring works on larger, single-output cones than rewriting. Driver
\id{Abc\_NtkRefactor} (\file{src/base/abci/abcRefactor.c}) computes a
\emph{reconvergence-driven cut} with \id{Abc\_NodeFindCut}
(\file{src/base/abci/abcReconv.c}), which greedily grows a leaf set that captures
reconvergent paths. Engine \id{Abc\_NodeRefactor} builds the cone truth table
(\id{Abc\_NodeConeTruth}), turns it into a fresh factored form with
\id{Kit\_TruthToGraph} (ISOP-based factoring), and measures the swap.

\begin{center}
\begin{tabular}{@{}lll@{}}
\toprule
\textbf{Parameter} & \textbf{Default} & \textbf{Role} \\
\midrule
\id{nNodeSizeMax} (\cmd{-N}) & 10 & max leaves of the refactored cut \\
\id{nConeSizeMax} (\cmd{-C}) & 16 & max support of the containing cone \\
\id{nMinSaved} (\cmd{-M})    & 1  & min nodes saved to accept \\
\bottomrule
\end{tabular}
\end{center}

The cut manager is started as
\id{Abc\_NtkManCutStart(nNodeSizeMax, nConeSizeMax, 2, 1000)}. A constant cone is
always accepted (returning \id{Dec\_GraphCreateConst0/1}). Otherwise the factored
form is rejected when \id{Dec\_GraphToNetworkCount} returns $-1$ (does not fit
under the required level) or when
$\text{nNodesSaved} - \text{nNodesAdded} < \text{nMinSaved}$. With \cmd{-z}
(\id{fUseZeros}) the threshold drops to $0$; with \cmd{-d} (\id{fUseDcs})
don't-cares from the containing cone are exploited.

\section{Resubstitution (\texorpdfstring{\cmd{resub}}{resub})}
\label{sec:opt-resub}

Resubstitution tries to re-express a node using signals (\emph{divisors}) that
already exist elsewhere in the window, deleting the node's MFFC. Driver
\id{Abc\_NtkResubstitute} (\file{src/base/abci/abcResub.c}) collects a window and
its divisors; engine \id{Abc\_ManResubEval} evaluates candidates. Divisor storage
is bounded by two compile-time caps:

\begin{center}
\begin{tabular}{@{}lll@{}}
\toprule
\textbf{Constant} & \textbf{Value} & \textbf{Meaning} \\
\midrule
\id{ABC\_RS\_DIV1\_MAX} & 150 & max single-node divisors considered \\
\id{ABC\_RS\_DIV2\_MAX} & 500 & max pair-wise divisors considered \\
\bottomrule
\end{tabular}
\end{center}

\paragraph{Fixed try order (first success wins).}
\id{Abc\_ManResubEval} attempts progressively larger reconnections and returns as
soon as one is valid, so cheap moves are always preferred:

\begin{center}
\begin{tabularx}{\linewidth}{@{}llX@{}}
\toprule
\textbf{Step} & \textbf{Routine} & \textbf{Result / gain} \\
\midrule
constant     & \id{Abc\_ManResubQuit}  & tie to const, gain \id{nMffc} \\
0-node       & \id{Abc\_ManResubDivs0} & equal divisor exists, gain \id{nMffc} \\
1-node       & \id{Abc\_ManResubDivs1} & one new gate, gain \id{nMffc}$-1$ \\
2-node       & \id{Abc\_ManResubDivs12}, \id{Abc\_ManResubDivs2} & gain \id{nMffc}$-2$ \\
3-node       & \id{Abc\_ManResubDivs3} & gain \id{nMffc}$-3$ \\
\bottomrule
\end{tabularx}
\end{center}

The search depth is gated by \cmd{-N} (\id{nNodesMax}, the max nodes to add,
range $0$--$3$, default $1$): with \id{nSteps==k} the loop stops after the
$k$-node stage. The window/cut size is \cmd{-K} (\id{nCutsMax}, range $4$--$16$,
default $8$); \cmd{-M} sets \id{nMinSaved} and \cmd{-F} enables observability
don't-cares over \id{nLevelsOdc} levels. The window resub helpers live in
\pkg{src/opt/res}.

\section{Don't-care / SAT-based (\texorpdfstring{\cmd{mfs}}{mfs})}
\label{sec:opt-mfs}

The strongest (and slowest) passes build a care-set/don't-care window around a
node and use a SAT solver plus interpolation to find a smaller function. Driver
\id{Abc\_NtkMfs} (\file{src/opt/mfs/mfsCore.c}) computes a window
(\id{Abc\_MfsComputeRoots}), formulates a miter, and calls
\id{Abc\_NtkMfsSolveSatResub} / \id{Abc\_NtkMfsInterplate}. Because it reasons
about the exact local don't-cares (not just structural cuts), it removes fanins
and nodes that structural methods miss --- at higher runtime cost. Defaults
(\id{Abc\_NtkMfsParsDefault}):

\begin{center}
\begin{tabular}{@{}lll@{}}
\toprule
\textbf{Field} & \textbf{Default} & \textbf{Meaning} \\
\midrule
\id{nWinTfoLevs}  & 2    & TFO levels included in the window \\
\id{nFanoutsMax}  & 30   & fanout cap when growing the window \\
\id{nDepthMax}    & 20   & max node level to process \\
\id{nWinMax}      & 300  & max window size (nodes) \\
\id{nGrowthLevel} & 0    & allowed level growth \\
\id{nBTLimit}     & 5000 & SAT conflict limit per call \\
\bottomrule
\end{tabular}
\end{center}

The mapped-network counterpart lives in \pkg{src/opt/sfm} (\id{Sfm\_ParSetDefault}
in \file{src/opt/sfm/sfmCore.c}) and uses the same window/SAT budgets
(\id{nTfoLevMax=2}, \id{nFanoutMax=30}, \id{nDepthMax=20}, \id{nWinSizeMax=300},
\id{nBTLimit=5000}).

\section{Cut enumeration utility}
\label{sec:opt-cut}

Several passes above rest on the $k$-feasible cut package in \pkg{src/opt/cut}.
\id{Cut\_ManStart} allocates a manager from a \id{Cut\_Params\_t}
(\id{nVarsMax}, \id{nKeepMax}, \id{fTruth}, \id{fFilter}, \ldots); cuts are merged
bottom-up in \id{cutNode.c} and truth tables maintained in \id{cutTruth.c}. The
cut size is bounded by \id{CUT\_SIZE\_MIN = 3} and \id{CUT\_SIZE\_MAX = 12}
(\file{src/opt/cut/cut.h}); the \cmd{cut} command exposes it via \cmd{-K}. The AIG
package embeds its own faster cut computation (\file{src/opt/dar/darCore.c},
\id{Dar\_ManComputeCuts}).

\section{Summary}
\label{sec:opt-summary}

\begin{table}[h]
\centering
\begin{tabularx}{\linewidth}{@{}llllX@{}}
\toprule
\textbf{Pass} & \textbf{Command} & \textbf{Driver} & \textbf{Engine} & \textbf{Key constants} \\
\midrule
Balance (AIG) & \cmd{balance} & \file{abcBalance.c} & \id{Abc\_NodeBalance\_rec} & supergate cap 10000 \\
Balance (GIA) & \cmd{\&b} & \file{giaBalAig.c} & \id{Gia\_ManBalanceGate} & supergate cap 50; lookback 8; slack $\le$100 \\
Rewrite & \cmd{rewrite} & \file{abcRewrite.c} & \id{Rwr\_NodeRewrite} & 4-input cuts; 222 classes; keep 250; fanout$>$1000 skip \\
DAR rewrite & \cmd{drw}/\cmd{dc2} & \file{darCore.c} & \id{Dar\_ManRewrite} & nCutsMax 8; nSubgMax 5; nMinSaved 1 \\
Refactor & \cmd{refactor} & \file{abcRefactor.c} & \id{Abc\_NodeRefactor} & node 10; cone 16; nMinSaved 1 \\
Resub & \cmd{resub} & \file{abcResub.c} & \id{Abc\_ManResubEval} & 150 divs; 500 pairs; K$=$8; N$=$0..3 \\
Don't-care & \cmd{mfs} & \file{mfsCore.c} & \id{Abc\_NtkMfs} & win 300; TFO 2; BT 5000 \\
\bottomrule
\end{tabularx}
\caption{Technology-independent optimization passes at a glance.}
\label{tab:opt-summary}
\end{table}

None of these passes is used alone. Scripts such as \cmd{resyn2},
\cmd{resyn2rs}, and \cmd{compress2} sequence balance/rewrite/refactor/resub into
robust recipes that iterate to a fixed point; these are cataloged in
Chapter~\ref{ch:recipes}.
